\chapter{How Unsupervised Learning Works}
\label{ch:unsupervised}

\chapterguide%
  {The shared foundations in Stages 1--2 and the supervised workflow in Stage 3A for comparison.}
  {define unsupervised objectives and evaluate structure using internal, external, and downstream evidence.}

Given features $X$ without labels, unsupervised learning finds groups, representations, or unusual cases.

\begin{mlfigure}
\begin{tikzpicture}[
  bx/.style={draw=MLBlue,fill=MLPaleBlue,rounded corners=2mm,text width=4.4cm,minimum height=1.2cm,align=center,font=\small\bfseries\color{MLNavy}},
  arr/.style={-{Stealth[length=2.2mm]},thick,draw=MLTeal}
]
\node[bx] (s) at (-4.3,0) {Supervised\\$X+y \longrightarrow$ predict $y$};
\node[bx] (u) at (4.3,0) {Unsupervised\\$X \longrightarrow$ discover structure};
\draw[arr] (s.east) -- node[above,font=\scriptsize\color{MLMuted}]{remove the supplied label} (u.west);
\end{tikzpicture}
\end{mlfigure}

\section{The unsupervised-learning map}

\begin{mltable}
\begin{tabularx}{\linewidth}{@{}>{\bfseries\leavevmode\color{MLNavy}}p{0.27\linewidth}p{0.27\linewidth}X@{}}
\toprule
\rowcolor{MLPaleBlue!92}
Question & First framing & Main output \\
\midrule
Do observations form groups? & Clustering & Cluster assignments, centers, hierarchy, or soft memberships \\
Can the data be represented with fewer dimensions? & Dimensionality reduction / representation & Components or lower-dimensional coordinates \\
Which observations are unusual? & Anomaly / novelty detection & Anomaly score or ranked suspicious cases \\
Where is the data concentrated? & Density estimation & Estimated probability density or mixture components \\
Which items occur together? & Association / pattern mining \textit{(beyond core)} & Frequent itemsets and rules \\
\bottomrule
\end{tabularx}
\end{mltable}

\begin{keyidea}
Evaluation follows the objective: variance for PCA, separation for clustering, reviewed-alert precision for anomalies.
\end{keyidea}

\section{What does fitting mean without labels?}

An unsupervised method still has an objective; it simply does not compare predictions with a supplied label during fitting.

\begin{mltable}
\begin{tabularx}{\linewidth}{@{}>{\bfseries\leavevmode\color{MLNavy}}p{0.24\linewidth}p{0.34\linewidth}X@{}}
\toprule
\rowcolor{MLPaleBlue!92}
Method family & What fitting tries to do & What is learned \\
\midrule
$K$-means & Reduce squared distance from points to assigned centers & Cluster centers and assignments \\
PCA & Preserve as much variance as possible in a lower-dimensional linear representation & Principal directions and projected coordinates \\
Density / mixture models & Assign high likelihood to the observed feature distribution under a chosen family & Distribution parameters and soft memberships \\
Isolation / local methods & Make unusual observations easy to isolate or locally sparse & Anomaly scores / rankings \\
Association mining \textit{(beyond core)} & Find item combinations that occur often enough to be interesting & Frequent itemsets and conditional rules \\
\bottomrule
\end{tabularx}
\end{mltable}

\section{How do we evaluate when there is no answer key?}

There is no universal ``unsupervised accuracy.'' Use evidence that matches the task.

\subsection{Internal evidence}
Assess compactness, separation, reconstruction, retained variance, and stability under resampling or perturbation.

\subsection{External evidence}
Compare with independent reference labels when available, using measures such as adjusted Rand index. Disagreement may reflect a different valid structure.

\subsection{Practical or downstream evidence}
Test usefulness: interpretable clusters, retained downstream signal, or anomalies that merit review.

\begin{pitfall}
Do not invent labels after fitting and treat them as ground truth. Known classes and visual groupings need not match the structure the method preserves.
\end{pitfall}

\section{The unsupervised workflow}

\begin{mlfigure}
\begin{tikzpicture}[
  bx/.style={draw=MLBlue,fill=MLPaleBlue,rounded corners=1.5mm,text width=2.0cm,inner xsep=2pt,minimum height=0.9cm,align=center,font=\scriptsize\bfseries\color{MLNavy}},
  arr/.style={-{Stealth[length=2mm]},thick,draw=MLTeal}
]
\node[bx] (a) at (-6.5,0) {1. Define the\\structure sought};
\node[bx] (b) at (-3.9,0) {2. Choose the\\unit and features};
\node[bx] (c) at (-1.3,0) {3. Prepare and\\scale appropriately};
\node[bx] (d) at (1.3,0) {4. Fit candidate\\structures};
\node[bx] (e) at (3.9,0) {5. Check stability\\and meaning};
\node[bx] (f) at (6.5,0) {6. Use on new or\\downstream data};
\draw[arr] (a)--(b); \draw[arr] (b)--(c); \draw[arr] (c)--(d); \draw[arr] (d)--(e); \draw[arr] (e)--(f);
\end{tikzpicture}
\end{mlfigure}

Describing a fixed dataset need not require a split. Claims about new data or downstream prediction do: fit on training rows and evaluate held-out cases, as in Labs~4--5.

\section{Scaling and representation matter even more here}

Scaling, weighting, encoding, and distance define geometry. Lab~4 standardizes age, income, and spending score before $K$-means.

\begin{keyidea}
There is no label to reveal a poor representation, so justify the scaling and distance choices before interpreting the result.
\end{keyidea}

\begin{quickcheck}
\begin{enumerate}
  \item Why is there no universal unsupervised-learning accuracy score?
  \item Give one internal, one external, and one downstream form of evidence that could support an unsupervised result.
  \item When does an unsupervised project need a held-out split even though no label is used during fitting?
\end{enumerate}
\end{quickcheck}
